\chapter{Ce que le hackathon attend}
\label{ch:attentes}

Ce chapitre rassemble tout ce qui a été demandé, écrit ou dit, sans interprétation. Il se termine
par une liste d'exigences numérotées : le chapitre~\ref{ch:couverture} vérifie que la solution
répond à chacune.

\section{La règle commune aux deux sujets}

Les slides de CGI posent quatre règles, \og deux sujets, une même règle : une IA frugale et open
source pour le domaine bancaire \fg.

\begin{center}\small
\begin{tabularx}{\textwidth}{@{}L{4cm} Y@{}}
\toprule
\textbf{Règle} & \textbf{Ce qu'elle impose} \\
\midrule
Open source uniquement & Licences ouvertes, aucune API de LLM externe, données en local. \\
Le plus petit modèle qui suffit & Le GPU est partagé : une base simple d'abord, un modèle plus gros
  seulement si le gain est \emph{mesuré}. \\
Serveur GPU sur site & Données et modèles restent en interne. \\
Mesurer sa frugalité & Taille du modèle, énergie (CodeCarbon) et latence, présentées dans le
  pitch. \\
\bottomrule
\end{tabularx}
\end{center}

Denis Barthou a précisé : la solution tourne uniquement sur les machines locales ; utiliser d'autres
IA pour \emph{développer} est autorisé ; dans la solution, on privilégie de petits LLM, \og des
algorithmes au lieu de LLM quand c'est possible \fg, des bases de données et d'autres types d'IA
(classifieurs). Le critère de frugalité retient le plus petit LLM \emph{en Go, quantisation
comprise}.

À l'oral, CGI a ajouté trois nuances qui orientent directement notre architecture :
\begin{itemize}
  \item la frugalité, c'est peu de tokens, \og mais surtout \fg{} peu d'énergie ;
  \item \og le plus petit modèle qui suffit \fg{} est une phrase \og à double tranchant \fg : une
        architecture peut appeler \textbf{un gros modèle à un instant précis, pour un besoin
        précis}, et de petits modèles le reste du temps --- \og c'est là que vous allez innover \fg ;
  \item le sujet 1 est \og de l'IA générative qu'on orchestre \fg, \og du chatbot, mais pas du
        chatbot traditionnel \fg.
\end{itemize}

\section{Le sujet 1}

\begin{cle}
\og Un client écrit : \emph{Je veux acheter un appartement.} Votre assistant doit comprendre, poser
les bonnes questions et lui proposer une piste adaptée à sa situation. \fg
\end{cle}

\paragraph{Objectifs.} Comprendre les questions des clients en langage naturel ; personnaliser
selon l'âge, les revenus et l'historique ; conseiller, c'est-à-dire recommander un produit bancaire
simple.

\paragraph{Démarche attendue, en quatre temps.} \emph{Cadrer} (quels besoins et quelles intentions
prioriser ?), \emph{Préparer} (rendre les données exploitables en français), \emph{Concevoir}
(quelle architecture, et pourquoi la plus sobre ?), \emph{Démontrer} (prouver la valeur et la
frugalité).

\paragraph{Livrables attendus.} Données préparées en français ; modèle de compréhension des
intentions ; chatbot relié à une base clients simulée ; moteur de conseils personnalisés ; mesures
de frugalité.

\paragraph{Données fournies.} Banking77 (intentions) ; Bank Marketing, \og enrichi de revenus et
d'historique simulés \fg{} ; le catalogue de produits est à construire à partir d'informations
publiques.

\paragraph{Les situations données en exemple.} Six sur les slides, trois de plus à l'oral.

\begin{center}\small
\begin{tabularx}{\textwidth}{@{}l L{4.3cm} L{3.2cm} Y@{}}
\toprule
\textbf{Thème} & \textbf{Message} & \textbf{Profil} & \textbf{Comportement attendu} \\
\midrule
Carte & \og J'ai perdu ma carte hier soir, je fais quoi ? \fg & Inès, 34 ans, cadre & Réagir vite,
  guider l'opposition, rassurer \\
Réclamation & \og Pourquoi j'ai payé 8\eur{} de frais ce mois-ci ? \fg & Marc, 52 ans, découverts
  fréquents & Expliquer les frais, orienter si besoin \\
Épargne & \og Je voudrais mettre de côté pour mes enfants. \fg & Sophie, 41 ans, deux enfants &
  Demander l'horizon, proposer une piste \\
Virement & \og Je dois envoyer 500\eur{} à mon fils. \fg & Paul, 68 ans, peu à l'aise en ligne &
  Guider pas à pas, avec des mots simples \\
Crédit & \og Je veux un crédit pour acheter une voiture. \fg & Karim, 29 ans, 1\,900\eur{} de
  revenus & Questionner, rester prudent sur l'endettement \\
Insatisfaction & \og Votre appli ne marche jamais ! \fg & Client mécontent, 3\textsuperscript{e}
  message du jour & Reconnaître la frustration, passer la main \\
\midrule
Logement & \og Je veux acheter un appartement. \fg & (énoncé principal) & Comprendre, questionner,
  proposer une piste \\
Prêt étudiant & \og Je veux faire un prêt étudiant, comment je fais ? \fg & (exemple oral) &
  Regard d'un jeune client \\
Jargon & Un mot ou un sigle bancaire non compris & (exemple oral) & Expliquer avec des mots simples \\
\bottomrule
\end{tabularx}
\end{center}

\paragraph{Précisions données à l'oral.} Les jeux de données sont \og indicatifs \fg{} : on peut en
croiser d'autres. La profondeur est à notre choix : \og un agent qui va jusqu'au bout d'un workflow
et qui fait lui-même l'action \fg{} (couper la carte selon la gravité), ou un assistant qui guide
seulement. Ils veulent \og notre regard \fg{} de jeunes clients sur le niveau de service attendu. Le
chatbot doit répondre \og immédiatement \fg. Les attentes seront précisées avec chaque équipe au
fil des jours.

\section{Ce que le jury regarde}

Jury de trois personnes : Adrien (CGI, aussi coach), Denis Barthou (ENSEIRB) et Jean-Michel (La
Banque Postale, à distance ; les équipes peuvent lui poser des questions). Il \og évalue votre
démarche autant que votre résultat \fg.

\begin{center}\small
\begin{tabularx}{\textwidth}{@{}c l Y@{}}
\toprule
\textbf{Poids} & \textbf{Critère} & \textbf{Formulation officielle} \\
\midrule
25\,\% & Pertinence métier & Utile pour un client ou un analyste \\
20\,\% & Qualité technique & Résultats mesurés et bien évalués \\
20\,\% & Frugalité & Taille, énergie, latence \\
15\,\% & Explicabilité & Chaque conseil ou alerte se justifie \\
10\,\% & Démonstration & Elle fonctionne et se comprend \\
10\,\% & Présentation & Claire et tenue en 15 minutes \\
\bottomrule
\end{tabularx}
\end{center}

\paragraph{Format.} 30 minutes par équipe (20 s'il y a plus de 7 groupes) : 15 minutes de
présentation avec \textbf{démonstration en direct sur le serveur fourni}, environ 10 minutes de
questions, 5 minutes de retour à chaud. À remettre avant : le code source déposé sur un dépôt Git
et le support de présentation.

\section{Ce que La Banque Postale valorise}

La directrice SI de La Banque Postale a présenté ses priorités IA dans cet ordre :
\begin{enumerate}
  \item \textbf{l'expérience client} : chatbot conversationnel sur ses contrats, \textbf{les
        réclamations} (\og un sujet prégnant \fg), et la personnalisation des offres à partir des
        avoirs du client --- c'est exactement le sujet 1 ;
  \item l'optimisation des processus (dossiers de crédit, anti-blanchiment, anti-fraude) ;
  \item l'innovation (outils pour les collaborateurs, incubateur).
\end{enumerate}
Elle a insisté sur trois valeurs : \textbf{l'accessibilité} (clients et collaborateurs malvoyants
ou malentendants : tout ce que construit la DSI doit être compatible), \textbf{l'inclusion
bancaire} (la banque a le devoir d'accueillir tous les clients, y compris les plus fragiles) et la
\textbf{responsabilité} (banque citoyenne, finance durable). CGI a mentionné la souveraineté et
la RSE parmi les enjeux de ses clients. Un client sur quatre en France est client de La Banque
Postale, accueilli dans 17\,000 points de contact.

\section{Les contraintes}

\begin{itemize}
  \item \textbf{Le temps.} Lundi soir, mardi, mercredi matin : un jour et demi de travail effectif.
  \item \textbf{Les machines.} Postes de la salle (RTX 3060 12 Go, 30 Go de RAM) : machine
        principale de développement. DeepEirb (2 $\times$ A6000 48 Go) : SSH et serveur
        d'inférence vLLM, à condition de \emph{demander à Denis Barthou} les modèles voulus.
        Jetson AGX Orin et OrangePi : disponibles mardi, pour le déploiement uniquement.
  \item \textbf{Les modèles déjà installés.} \f{gemma4:e2b} (7,2 Go) et \f{ministral-3:14b}
        (9,1 Go) sous Ollama, contexte réglé à 16\,384 tokens. Télécharger d'autres modèles est
        possible, \og pas en wifi \fg.
  \item \textbf{L'équipe.} Trois personnes : Yahya (IA), Sohayb (génie logiciel), Younes (calcul
        intensif).
\end{itemize}

\section{Ce qu'un organisateur nous a conseillé}

Lundi vers 17h30, un organisateur du hackathon a pris le temps d'échanger avec nous
(transcription dans \f{infos/transcription-3-echange-organisateur.txt}). Ses conseils précisent
la règle commune :

\begin{itemize}
  \item \textbf{\og Pour gagner, c'est moins de LLM, plus d'algo \fg}, ou des algorithmes différents,
        ou des modèles d'IA plus spécialisés --- \og même pour le premier \fg{} sujet.
  \item \textbf{En entrée, pas besoin de LLM} pour savoir de quoi parle le client : un algorithme de
        catégorisation est plus rapide. Une simple recherche de \textbf{mots-clés}, presque gratuite,
        trouve le sujet dans environ 30\,\% des cas (\og ça fait 30\,\% de gagné \fg), et peut
        donner \textbf{un degré de confiance} : un seul mot-clé trouvé dans une longue phrase, ce
        n'est pas sûr. Les autres cas passent à un outil plus lourd.
  \item \textbf{En sortie, un LLM}, parce qu'il faut parler français à un humain ; mais il faut
        \textbf{minimiser les appels} et ne l'appeler \textbf{qu'à la fin}.
  \item \textbf{Choisir sa taille en descendant} : commencer avec un modèle qui parle correctement
        (\og un 15B ou un 30B \fg), puis descendre vers un 4B, un 2B (\og même le Gemma 4 que vous
        avez \fg), voire un 1B.
  \item \textbf{Se demander, à chaque endroit, si un LLM est nécessaire} : jamais, ou seulement pour
        une partie des cas. \og Seulement 30\,\% des cas avec un LLM, c'est intéressant \fg{} : on
        prend le vélo pour aller chercher le pain, la voiture pour le week-end.
  \item \textbf{Un seul sujet}, bien fait. Puis, \textbf{une fois la solution sûre}, essayer des
        idées nouvelles, même exotiques ; utiliser l'IA pour coder et discuter des idées ; et
        \og se faire plaisir \fg.
\end{itemize}

\section{Les exigences, numérotées}
\label{sec:exigences}

\begin{center}\small
\begin{longtable}{@{}l L{8.6cm} L{4.4cm}@{}}
\toprule
\textbf{N\textsuperscript{o}} & \textbf{Exigence} & \textbf{Source} \\
\midrule
\endhead
E1 & Comprendre les demandes en langage naturel, en français & Objectif 01, livrable \\
E2 & Fournir des données préparées en français & Livrable \\
E3 & Poser les bonnes questions avant de proposer & Énoncé \\
E4 & Personnaliser selon l'âge, les revenus et l'historique & Objectif 02 \\
E5 & Recommander un produit bancaire simple & Objectif 03, livrable \\
E6 & Relier le chatbot à une base clients simulée & Livrable \\
E7 & Construire le catalogue à partir d'informations publiques & Données \\
E8 & Open source, aucune API externe, données locales & Règle commune \\
E9 & Le plus petit modèle qui suffit ; un plus gros seulement si le gain est mesuré & Règle commune \\
E10 & Mesurer taille, énergie (CodeCarbon) et latence & Règle commune, livrable \\
E11 & Justifier chaque conseil & Grille (15\,\%) \\
E12 & Répondre immédiatement & Oral \\
E13 & Choisir et assumer la profondeur : guider ou agir & Oral \\
E14 & Traiter les six situations avec le comportement attendu & Slide 3 \\
E15 & Expliquer le jargon ; apporter le regard d'un jeune client & Oral \\
E16 & Démo en direct sur le serveur fourni ; code sur Git ; support remis & Soutenance \\
E17 & Respecter accessibilité et inclusion (souhaitable) & Présentation de La Banque Postale \\
E18 & En entrée, des algorithmes ; un seul appel au LLM, à la fin ; mesurer la part des cas sans
  LLM & Échange avec un organisateur \\
E19 & D'abord une solution sûre, ensuite des idées nouvelles & Échange avec un organisateur \\
\bottomrule
\end{longtable}
\end{center}
